% from ega4-18-fr.tex
\oldpage[IV-4]{213}

Enfin, lorsque $A$ est noethérien, que les $f_i$ appartiennent au radical de $A$, que la suite $\mathbf f$ est $\operatorname{gr}_{\mathfrak R}^*(A)$-régulière et que $M$ est un $A$-module de type fini, les conditions $a)$, $b)$, $c)$, $d)$ sont équivalentes.
\end{theorem}

L'implication $a)\Rightarrow c)$ est immédiate ($0$, 15.1.13). Si $c)$ est vérifiée, il résulte de (19.5.5) que $\psi_M$ est bijectif; en outre, $\operatorname{gr}_{\mathfrak R}^*(M)\otimes_A(A/\mathfrak L)=\operatorname{gr}_{\mathfrak R}^*(M)\otimes_{A/\mathfrak R}(A/\mathfrak L)$ est un $(A/\mathfrak L)$-module plat, donc il en est de même de $(\operatorname{gr}_{\mathfrak R}^*(M)\otimes_A(A/\mathfrak L))[T_1,\ldots,T_n]$, et par suite de $\operatorname{gr}_{\mathfrak L}^*(M)$ puisque $\psi_M$ est un $(A/\mathfrak L)$-isomorphisme; donc $c)$ entraîne $b)$. Le fait que $a)$ implique $b)$ montre aussitôt que $a)$ implique aussi $d)$.

Lorsque tout quotient d'un $\operatorname{gr}_{\mathfrak R}^p(M)$ est idéalement séparé pour $\mathfrak J$, il résulte de (19.5.5) que la condition $b)$ entraîne que $\mathbf f$ est $(\operatorname{gr}_{\mathfrak R}^*(M))$-régulière; en outre, comme $(\operatorname{gr}_{\mathfrak R}^*(M)\otimes_{A/\mathfrak R}(A/\mathfrak L))[T_1,\ldots,T_n]$, isomorphe à $\operatorname{gr}_{\mathfrak L}^*(M)$, est alors un $(A/\mathfrak L)$-module plat, il en est de même de $\operatorname{gr}_{\mathfrak R}^*(M)\otimes_{A/\mathfrak R}(A/\mathfrak L)$, qui en est un facteur direct; cela prouve donc que $b)$ entraîne alors $c)$. D'autre part, sous la même hypothèse, $c)$ entraîne $a)$ en vertu de ($0$, 15.1.21) (où on peut remplacer l'hypothèse noethérienne par l'hypothèse que $M$ est idéalement séparé pour $\mathfrak J$, en vertu de ($0_{\mathrm{III}}$, 10.2.1)).

Il reste donc à prouver que lorsque $A$ est noethérien, $M$ de type fini et que les $f_i$ appartiennent au radical de $A$, $d)$ entraîne $a)$.

Il existe par hypothèse un $A$-module \emph{libre de type fini} $N$ et une suite exacte $0\to R\to N\to M\to0$. Il suffit de prouver que $d)$ entraîne $b)$, puisque tout quotient d'un sous-module d'un $\operatorname{gr}_{\mathfrak R}^p(M)$ est alors idéalement séparé pour $\mathfrak J$, en vertu du fait que $\operatorname{gr}_{\mathfrak R}^p(M)$ est un $A$-module de type fini, que $\mathfrak J$ est contenu dans le radical de $A$, et que $A$ est noethérien (Bourbaki, \emph{Alg. comm.}, chap. III, § 5, n\textsuperscript{o} 1); en d'autres termes, il s'agit de voir que $\psi_M$ est \emph{bijectif}.

Comme par hypothèse la suite $\mathbf f$ est $\operatorname{gr}_{\mathfrak R}^*(A)$-régulière, elle est aussi $\operatorname{gr}_{\mathfrak R}^*(N)$-régulière; on peut par suite appliquer (19.6.3), car tout module quotient d'un sous-module d'un $\operatorname{gr}_{\mathfrak R}^k(M)$ est alors séparé pour la topologie $\mathfrak J$-préadique puisque les $f_i$ appartiennent au radical de $A$ par hypothèse. Le diagramme commutatif (19.6.2.1) s'écrit donc ici
\begin{equation}
\xymatrix{
  & 0 \ar[d] \\
  (\operatorname{gr}_{\mathfrak R}^*(R,N)\otimes_{A/\mathfrak R}(A/\mathfrak L))[T_1,\ldots,T_n] \ar[r]^-{\psi_R'} \ar[d] & \operatorname{gr}_{\mathfrak L}^*(R,N) \ar[d] \\
  (\operatorname{gr}_{\mathfrak R}^*(N)\otimes_{A/\mathfrak R}(A/\mathfrak L))[T_1,\ldots,T_n] \ar[r]^-{\psi_N} \ar[d] & \operatorname{gr}_{\mathfrak L}^*(N) \ar[d] \\
  (\operatorname{gr}_{\mathfrak R}^*(M)\otimes_{A/\mathfrak R}(A/\mathfrak L))[T_1,\ldots,T_n] \ar[r]^-{\psi_M} \ar[d] & \operatorname{gr}_{\mathfrak L}^*(M) \ar[d] \\
  0 & 0
}
\tag{19.7.1.1}\label{IV.19.7.1.1}
\end{equation}
où $\operatorname{gr}_{\mathfrak R}^*(R,N)$ (resp. $\operatorname{gr}_{\mathfrak L}^*(R,N)$) est le gradué associé à $R$ pour la filtration des $R\cap\mathfrak R^mN$ (resp. $R\cap\mathfrak L^mN$); rappelons que dans ce diagramme les colonnes sont exactes et que $\psi_N$

